\subsection{Explicit rules, recursion, and long-term behavior}

\begin{definitionbox}[Sequence]
A real sequence is a function
\[
a:\mathbb N\to\mathbb R.
\]
Its value at the input $n$ is usually written $a_n$ rather than $a(n)$.
\end{definitionbox}

A sequence may be given explicitly, for example
\[
a_n=\frac{1}{n},
\]
or recursively, for example
\[
a_1=1,
\qquad
a_{n+1}=\frac12a_n.
\]

\begin{examplebox}[Reading an explicit sequence]
For $a_n=1/n$, the first terms are
\[
1,\frac12,\frac13,\frac14,\ldots
\]
They remain positive and get closer to zero as $n$ grows.
\end{examplebox}

\begin{examplebox}[Reading a recursive sequence]
If $a_1=8$ and $a_{n+1}=a_n/2$, then
\[
8,4,2,1,\frac12,\ldots
\]
Each term is determined from the previous one.
\end{examplebox}

\begin{interpretationbox}[Discrete input]
A sequence does not fill an interval of inputs. Its graph consists of separated points indexed by natural numbers. This is why the long-term question is phrased as $n\to\infty$.
\end{interpretationbox}
